% OWNER: theory-b
% STATUS: outline
% SOURCES: notes/SPIS_TRESCI_ROBOCZY.md § Rozdz. 5
\chapter{Sztuczna inteligencja generatywna i personalizacja w HCI}
\label{ch05:top}
\ChapterStatus{outline}{theory-b}

\section{Generatywne modele obrazu a architektura wnętrz}
\label{ch05:sec:genai}
\Todo{Możliwości i ograniczenia modeli (FLUX / Gemini) w domenie wnętrz.}

\section{Od promptu do koncepcji}
\label{ch05:sec:prompt}
\Todo{Kontrola stylu, wierność briefowi, artefakty generacji.}

\section{Explainable AI w systemach rekomendacyjnych}
\label{ch05:sec:xai}
\Todo{Przejrzystość mapowań preferencje → parametry → prompt.}

\section{Research through Design}
\label{ch05:sec:rtd}
\Todo{Artefakt jako metoda badawcza.}
\NeedsCite{zimmerman2007rtd}

\section{Gamifikacja instrumentów psychologicznych}
\label{ch05:sec:gamifikacja}
\Todo{Tinder, mood grid, awatar IDA — uzasadnienie HCI.}

\section{Etyka badań online i RODO}
\label{ch05:sec:etyka}
\Todo{Zgoda, anonimizacja (\texttt{user\_hash}), wycofanie; załącznik~\ref{app:f:top}.}

\section{Pozycjonowanie IDA względem narzędzi AI do wnętrz}
\label{ch05:sec:pozycjonowanie}
\Todo{Tabela porównawcza: konsumenckie AI vs artefakt badawczy.}
